\section{Deleting two primes}\label{sec:local}
The eight reduced residue classes modulo $24$ are
$1,5,7,11,13,17,19,23$. Plichta's prime cross draws these classes
as spokes of a wheel. Grant and Ghannam use the same picture
\cite{plichta,grant}. A prime greater than $3$ lies on one of these
spokes, as does a composite such as $91$. For the von Mangoldt function
$\Lambda$, which takes the value $\log p$ on $p^k$, $k\geq1$, and is
zero elsewhere, restricting to these eight classes has a precise effect.
It removes only the powers of $2$ and $3$.

We consider a weight used by Alpöge \cite[p.~3, item 5]{alpoge}.
For $x\geq1$ real and $m\geq1$ an integer, define
\begin{align}
 S_m(x)&=\sum_{\substack{n\leq x\\(n,m)=1}}\Lambda(n)(x/n-1),
 &c_m&=\sum_{p\mid m}\frac{\log p}{p-1},\notag\\
 E_m(x)&=S_m(x)-x\log x+(1+\gamma+c_m)x. \label{eq:normalization}
\end{align}
Here $\gamma$ is Euler's constant. Write $E=E_1$, $\ell_p=\log p$
and $a_p=\ell_p/(p-1)$. For $0\leq v\leq1$ define
\[
 f_p(v)=a_p e^{\ell_p v}-\ell_p v.
\]
Since $f_p(0)=f_p(1)=a_p$, this gives a continuous function on the
circle $\T=\RR/\mathbb Z$. All measures on tori below are normalized
Lebesgue measures. We use natural logarithms throughout.

\begin{prop}[Exact correction]\label{prop:local}
For every $m\geq1$ and $x\geq1$,
\begin{equation}
 E_m(x)=E(x)+\omega(m)\log x+
 \sum_{p\mid m}f_p\!\left(\frac{\log x}{\ell_p}\right),
 \label{eq:correction}
\end{equation}
where $\omega(m)$ is the number of distinct prime divisors of $m$.
\end{prop}
\begin{proof}
For a prime $p\mid m$ put $K=\lfloor\log x/\ell_p\rfloor$ and
$u=x/p^K$. The contribution of its deleted powers to $S_1-S_m$ is
\[
 \ell_p\sum_{k=1}^K(xp^{-k}-1)
 =a_p(x-u)-K\ell_p.
\]
The normalization restores $a_px$, and
$K\ell_p=\log x-\log u$.
The remaining term $a_pu-\log u$ is exactly
$f_p(\log x/\ell_p)$. Distinct primes have disjoint sets of positive
powers. At $x=p^K$ the last summand vanishes, so the formula includes
every endpoint.
\end{proof}

For the rest of the distributional analysis fix distinct primes $p,q$.
Set $\omega=(1/\ell_p,1/\ell_q)$,
$F(v,w)=f_p(v)+f_q(w)$, and $R_{p,q}(x)=F(\omega\log x)$.
Write $R_{24}=R_{2,3}$. Logarithmic oscillations in sums over geometric
progressions are classical in Mellin analysis \cite{fgd}.
The one-base distributional behavior of $(\log_b n)\bmod 1$ under
integer sampling has been studied quantitatively by Xu and Berger
\cite{xu,berger}.
Here the two explicit local functions determine the support, both
endpoint constants, and the image of the joint phase distribution.
All distributional statements below are independent of hypotheses on
zeros of $L$-functions.
